\section{Scalar jet bounds}
\label{sec:scalar}

The next result compares section orders and curve degrees with the
volume of an ample rational divisor. Its hypothesis on equivariant
fibrations is essential to both estimates. In the index argument,
Lemma~\ref{lem:index-general-type-fibres} establishes this
hypothesis on the varieties to which we apply the result.

We use the normalized order invariant of Definition~\ref{def:orders}.
The following estimates apply to rational polarizations of arbitrary
denominator.  All orders and section spaces are taken in sufficiently
divisible degrees; taking powers shows that imposing any further fixed
divisibility does not change the normalized supremum.

\begin{theorem}[Scalar bounds]\label{thm:scalar-bounds}
There are constants $C,c>0$ with the following property.  Let $V$ be a
normal projective canonical fourfold with $K_V\sim_{\Q}0$, and let a finite
group $G$ act on $V$.  Suppose that every $G$-equivariant rational fibration
with geometrically integral generic fibre and with both base and fibre of
positive dimension has geometric generic fibre of general type on a
smooth projective resolution.  It is enough to impose this condition on
the rational quotients supplied by Lemma~\ref{lem:generic-chain}.
If $P$ is an ample rational Cartier divisor satisfying
$g^*P\sim_{\Q}P$ for every $g\in G$, then
\begin{equation}\label{eq:scalar-bounds}
 \gamma(P;V)\le C(P^4)^{1/4},
 \qquad
 \varepsilon(P;x)\ge c(P^4)^{1/4}
\end{equation}
at a very general smooth point $x\in V$, where
\[
 \varepsilon(P;x)=
 \inf_{\substack{x\in D\subset V\\D\text{ an integral curve}}}
 \frac{P\cdot D}{\mult_xD}.
\]
The constants are independent of $V$, $G$, and $P$.
\end{theorem}

The upper estimate is the main geometric step. If a section can have
large normalized order on a polarization of very small volume, its
moving base locus produces smaller-dimensional covering members.
Numerical subadjunction bounds their canonical degrees. Repeatedly
passing to their Iitaka fibres will produce members of Kodaira dimension
zero with very small section orders. If their chains fill $V$, the
chain-order bound contradicts the initial large section order. Otherwise
they give a proper general-type leaf covered by Kodaira-dimension-zero
members, contradicting Lemma~\ref{lem:scalar-covering}. Once the upper
bound is proved, a thickening of a small-degree chain leaf gives the
lower curve bound.

We first establish the local numerical input.  For rational Cartier
classes $A,B$, we write $A\preceq B$ when $B-A$ is rationally linearly
equivalent to an effective rational divisor.  A restriction of such a
comparison will always be made to a member not contained in that effective
divisor.

\subsection{Numerical subadjunction and moving base components}

\begin{lemma}[Numerical subadjunction at a generic lc centre]
\label{lem:scalar-subadjunction}
Let $Z$ be a projective $\Q$-factorial klt variety, let $\Theta\ge0$ be a
rational divisor, and let $S\subset Z$ be an integral subvariety of
dimension $d>0$.  Suppose that $(Z,\Theta)$ is lc over a neighbourhood of
the generic point of $S$, and that $S$ is an lc centre there.  For a nef
rational Cartier divisor $A$ on $Z$ and a smooth projective resolution
$S^*\to S$, one has
\begin{equation}\label{eq:scalar-subadjunction}
 K_{S^*}\cdot(A|_{S^*})^{d-1}
 \le (K_Z+\Theta)|_S\cdot(A|_S)^{d-1}.
\end{equation}
Here the right-hand side means the intersection of the restricted
rational line bundle with the cycle $S$.
\end{lemma}

\begin{proof}
The arbitrary-boundary dlt blow-up theorem
\cite[Theorem~2.10]{FujinoHashizume} gives a projective birational morphism
$p:Q\to Z$, with $Q$ $\Q$-factorial, such that
\begin{equation}\label{eq:scalar-dlt-excess}
 K_Q+B_Q+E=p^*(K_Z+\Theta),
 \qquad (Q,B_Q)\text{ dlt},\qquad E\ge0,
\end{equation}
and $\Supp E$ lies above the non-lc locus of $(Z,\Theta)$.
The theorem allows coefficients of $\Theta$ greater than one: $B_Q$ is
the truncated strict transform together with the reduced exceptional
divisor, and the remaining crepant boundary is the effective excess $E$.

Choose a dlt stratum $R$ minimal by inclusion among the strata mapping
onto $S$.  Such a stratum exists.  Indeed, over the lc neighbourhood of
the generic point of $S$, the morphism in
\eqref{eq:scalar-dlt-excess} is crepant, and the centre on $Q$ of an lc
place with centre $S$ is a dlt stratum mapping onto $S$.
Dlt adjunction, in the form of the lc-centre and different theorem
\cite[4.16]{Kollar2013}, gives a normal $R$ and an effective
different.  Adding the restriction of $E$ gives an effective rational
divisor $\Delta_R$ with
\[
 K_R+\Delta_R\sim_{\Q}(K_Z+\Theta)|_R.
\]
The restriction of $E$ is defined and effective: $E$ is rational Cartier
and misses the generic fibre of $R\to S$.  Moreover $(R,\Delta_R)$ is
klt over the generic point of $S$.  Any lc centre of the dlt different
dominating $S$ would give a smaller dlt stratum of $Q$ dominating $S$,
contrary to the choice of $R$; the added excess is zero over that generic
point.

Let $S^\nu$ be the normalization of $S$ and factor the map through its
normal Stein base:
\[
 R\xrightarrow{f}S'\xrightarrow{h}S^\nu\longrightarrow S,
 \qquad f_*\OO_R=\OO_{S'},\qquad h\text{ finite}.
\]
Write $D$ for the pullback of $(K_Z+\Theta)|_S$ to $S^\nu$.
The morphism $f:(R,\Delta_R)\to S'$ satisfies the definition of a
klt-trivial fibration in \cite[Definition~3.1]{FujinoGongyoModuli}.
It is projective, generically subklt, and its adjoint is rationally
linearly equivalent to $f^*h^*D$.  The rank-one condition also holds:
on a resolution over the generic base point, the rounded discrepancy
divisor is effective and exceptional over $R$.  Allowing poles along
such a divisor does not enlarge regular functions on the normal variety
$R$, and connectedness of the fibres then gives rank one.

Ambro's canonical bundle formula, in the form recalled in
\cite[Theorem~3.5]{FujinoGongyoModuli}, therefore gives
\begin{equation}\label{eq:scalar-cbf}
 h^*D\sim_{\Q}K_{S'}+B_{S'}+M_{S'},
\end{equation}
where $M_{S'}$ is the trace of b-nef rational data.  The discriminant
$B_{S'}$ is effective.  To see this even where the total pair is not lc,
let $P$ be a prime of $S'$.  Over its generic point, $f^*P$ has a divisor
of multiplicity at least one and $\Delta_R$ has nonnegative coefficient
on that divisor.  The corresponding threshold is therefore at most
one.  It can be negative, but its discriminant coefficient, one minus
the threshold, remains nonnegative.

Let $A'$ be the pullback to $S'$ of $A|_{S^\nu}$.  Intersect
\eqref{eq:scalar-cbf} with $(A')^{d-1}$.  Both the effective discriminant
and the b-nef moduli trace contribute nonnegatively.  For the latter
assertion, compute on a model where the moduli divisor is nef and apply
the projection formula.  If $e=\deg h$, codimension-one ramification
for the finite map of normal varieties gives
\[
 K_{S'}\cdot(A')^{d-1}
 \ge e\,K_{S^\nu}\cdot(A|_{S^\nu})^{d-1}.
\]
These are intersections of Weil divisors with Cartier powers, so the
calculation can be made on the smooth codimension-one loci.  Finally,
the pushforward of $K_{S^*}$ is $K_{S^\nu}$, and its exceptional
divisors pair to zero with the pulled-back power of $A$.  Dividing by
$e$ proves \eqref{eq:scalar-subadjunction}.
\end{proof}

The moving-base-component strategy below adapts the jets construction in
\cite[Sections~7--8]{OpenAIAbundance}.  This method credit is distinct
from the qualitative nonvanishing input in the introduction; the public
jet-base-locus method and the present estimates are specified in the proof.

\begin{lemma}[Tracking a moving base component]\label{lem:scalar-tracking}
Let $Z$ be a projective $\Q$-factorial klt variety of dimension $n\ge2$,
and let $N$ be a big nef semiample rational Cartier divisor.  Suppose
there are $n+1$ positive rational numbers, with consecutive spacing
$\delta>0$, all strictly between $(N^n)^{1/n}$ and $\gamma(N;Z)$.
There is a covering family of integral subvarieties $S\subset Z$ of
some dimension $0<d<n$ such that, on a smooth projective resolution
$S^*$ of its geometric generic member,
\begin{align}
 N^d\cdot S&\le (2/\delta)^{n-d}N^n,
       \label{eq:scalar-track-degree}\\
 K_{S^*}\cdot(N|_{S^*})^{d-1}
 &\le \bigl(K_Z+(2n/\delta)N\bigr)|_S
             \cdot(N|_S)^{d-1}.
       \label{eq:scalar-track-canonical}
\end{align}
\end{lemma}

\begin{proof}
The moving-parameter argument follows the jet-base-locus method of
\cite[Proposition~2.3]{EinKuchleLazarsfeld1995}. We keep the
subadjunction and degree estimates explicit for the rational
polarizations used here.
Choose a sufficiently divisible positive integer $l$, clearing the
denominators of $N$ and of the levels.  It can be chosen arbitrarily
large so that every jet space under consideration is nonzero: first
choose a section witnessing an order above the highest level, and
then take powers.  At a variable very general smooth point $x$, write
\[
 V_\tau(x)=
 H^0\bigl(Z,\OO_Z(lN)\otimes\mathfrak m_x^{l\tau}\bigr).
\]
These spaces are the fibres of generic jet kernels, and thus have
rationally varying bases after restricting the parameter space.

At each level the base locus through $x$ is proper and has no isolated
component at $x$.  For the second assertion, if its local base ideal
were primary to $\mathfrak m_x$, then $n$ general members would have
local intersection multiplicity at least $(l\tau)^n$.  Successive
proper cuts, discarding other base components before each cut, would
give
\[
 l^nN^n\ge(l\tau)^n,
\]
contrary to $\tau>(N^n)^{1/n}$.  This degree argument is also detailed
below for the positive-dimensional component that we use.

The base loci are nested as the level increases.  Choose successively
an irreducible component through $x$, each contained in the next.
Their dimensions are in $\{1,\ldots,n-1\}$, so two consecutive
components agree.  Denote their common reduced component by $S$,
and spread it over a parameter space containing the varying mark $x$.
A finite extension of the parameter field may be used to name a
geometric component.  The resulting incidence dominates $Z$ because
it contains the graph of the mark.

Let $\tau$ and $\tau+\delta$ be the two levels.  We claim that the
base ideal of $V_{\tau+\delta}(x)$ has order at least $l\delta/2$ at
the generic point of $S$.  Differentiate a rationally varying basis
of this jet space in parameter directions, keeping the spatial
variable and a fixed basis of $H^0(Z,lN)$ constant.  A derivative of
length $r<l\delta$ has order at least $l(\tau+\delta)-r>l\tau$
at the varying mark, so it belongs to the lower jet space and
vanishes along the incidence of $S$.

Here these derivative vanishings detect the order of the original
section along the incidence.  At its generic point, the ambient
product is regular and both projections are dominant.  Since the
incidence dominates the spatial factor in characteristic zero,
parameter derivations span the dual of its conormal space.  The
associated graded ring of the regular local ring is the symmetric
algebra of that conormal space.  A nonzero initial form of degree
$r$ is detected by some iterated parameter derivative of length $r$,
by the characteristic-zero symmetric-power pairing.  Thus all the
vanishings just obtained force order at least $l\delta$, and in
particular the asserted weaker bound.  This is the same local Taylor
argument used in Lemma~\ref{lem:chain-order}.  Passing to the
geometric generic parameter fibre preserves it.

Put $a=n-d$.  The local ring of $Z$ at the generic point of $S$ is
regular of dimension $a$: a sweeping member meets the smooth open of
$Z$.  Since $S$ is a component of the higher base locus, its base
ideal there is primary to the maximal ideal $\mathfrak m$, and it is
contained in $\mathfrak m^{\lceil l\delta/2\rceil}$.  General choices
of $a$ members form a parameter ideal.  Their intersection length,
equal to its multiplicity in this regular local ring, is at least
$(l\delta/2)^a$.  Globally, take the cuts successively and discard
components already contained in the base locus before the next cut.
No discarded component contains the generic point of $S$, since $S$
is itself an irreducible component of that base locus.  The cuts are
therefore proper at the generic point being tracked.  Nefness of $N$
makes all discarded degree contributions nonnegative, and gives
\[
 (l\delta/2)^a(N^d\cdot S)\le l^aN^n.
\]
This proves \eqref{eq:scalar-track-degree}, as well as the isolated
base-point calculation used above.

Let $t$ be the local lc threshold of this base ideal over the generic
point of $S$.  A log resolution shows that $t$ is positive and rational.
Blowing up the smooth transverse centre, whose log discrepancy is
$a$, gives
\[
 t\le\frac{2a}{l\delta}.
\]
For $a=1$ the same inequality follows by testing $S$ itself.  Average
sufficiently many general members of the higher jet space and
multiply this average by $t$.  On a log resolution of the ideal the
free parts are general and their coefficients can be made arbitrarily
small.  This constructs an effective rational divisor
\[
 \Theta\sim_{\Q}bN,
 \qquad b=lt\le\frac{2a}{\delta},
\]
such that $(Z,\Theta)$ is lc near the generic point of $S$ and $S$ is
an lc centre there.  Indeed a place computing the threshold has
positive base-ideal order; the primary condition makes its centre
equal to $S$ after localization.  Lemma~\ref{lem:scalar-subadjunction}
with $A=N$ now gives \eqref{eq:scalar-track-canonical}, since $a\le n$.
\end{proof}

We will also need the following elementary consequence of a covering
family.

\begin{lemma}[Canonical dimension of a covering member]
\label{lem:scalar-covering}
Let $X$ be a projective variety birational to a smooth projective
variety with a nonzero pluricanonical section.  A smooth projective
resolution of the geometric generic member of a covering family of
integral subvarieties of $X$ has nonnegative Kodaira dimension.  If
$X$ is of general type, that resolution is of general type.
\end{lemma}

\begin{proof}
Pass to a smooth projective model of $X$ and to strict transforms of
the moving members.  Shrink the parameter space so that the incidence
has fibres of pure constant dimension.  If the incidence has dimension
larger than $\dim X$, intersect the parameter space in a projective
closure with general hyperplanes, in number equal to that excess.
Over the generic point of $X$ the incidence fibre embeds into the
parameter space, so a general such section meets its top-dimensional
part and makes the evaluation generically finite.  A component of
the cut incidence dominating $X$ dominates a component of the parameter
slice, by dimension.

Choose the cutting hyperplanes with independent generic coefficients
over the original field of definition.  A generic point of the selected
parameter slice remains generic in the original parameter space over
that field: every original proper closed subset has smaller dimension
after the same general cuts.  Thus the argument tests the original
geometric generic member, rather than a special subfamily.

Resolve and compactify the cut incidence and its evaluation to $X$.
The resulting generically finite morphism between smooth projective
varieties satisfies
\[
 K_I=q^*K_X+R,\qquad R\ge0.
\]
Restrict to a smooth geometric generic fibre of the parameter map.
It is birational to the corresponding member, and adjunction identifies
the restriction of $K_I$ with its canonical divisor.  A nonzero
pluricanonical section of $X$ pulls back to a section not identically
zero on this generic fibre, by dominance.  This proves nonnegative
Kodaira dimension.  If $K_X$ is big, restrict a rational decomposition
of a multiple into an ample class and an effective divisor.  The
ample pullback is big on the generic parameter fibre, since evaluation
is generically finite on its image.  Its canonical class is then big.
\end{proof}

\subsection{Small adjoint volumes on Iitaka fibres}

The tracking lemma produces a moving member of dimension at most three
whose polarization has small volume and whose canonical degree is bounded
by that volume.  The next lemma isolates the lower-dimensional argument:
its Iitaka fibres have Kodaira dimension zero, and adding the restricted
polarization to their canonical divisor still gives small volume.

\begin{lemma}\label{lem:scalar-iitaka-volume}
Fix $1\le d\le3$ and a positive rational number $C'$.  Let
$(W_i,L_i)$ be a sequence of smooth projective $d$-folds with
$\kappa(W_i)\ge0$ and big, nef, semiample rational Cartier divisors $L_i$.
Suppose that
\begin{equation}\label{eq:scalar-fibre-hypotheses}
 L_i^d\longrightarrow0,
 \qquad K_{W_i}\cdot L_i^{d-1}\le C'L_i^d.
\end{equation}
After passing to a subsequence and resolving the Iitaka fibrations,
their smooth geometric generic fibres $U_i$ have a fixed positive
dimension, satisfy $\kappa(U_i)=0$, and obey
\[
 \vol(K_{U_i}+L_i|_{U_i})\longrightarrow0.
\]
Here $L_i$ denotes its pullback to the resolved model; its restriction
to $U_i$ is big, nef, and semiample.  When $\kappa(W_i)=0$, the fibre
$U_i$ means $W_i$ itself.
\end{lemma}

\begin{proof}
Suppress the sequence index.  We first bound the volume of every fixed
positive adjoint combination on $W$, then compare it with the volume on
an Iitaka fibre.
For every fixed positive rational number $a$, we claim that
\begin{equation}\label{eq:scalar-small-adjoint-volume}
 \vol(aK_W+L)\longrightarrow0.
\end{equation}
Choose a general divided member of a sufficiently divisible free
multiple of $L/a$.  It gives an effective klt boundary
$\Gamma\sim_{\Q}L/a$.  This boundary is big, and so is
$K_W+\Gamma$, since $\kappa(W)\ge0$.
The klt log-general-type model and finite-generation theorem
\cite[Theorem~1.2]{BCHM2010} gives a $\Q$-factorial klt log terminal model.
Its nef big adjoint is semiample by the klt basepoint-free theorem
\cite[Theorem~3.3]{KollarMori1998}: apply it to a Cartier multiple of
the adjoint, whose sufficiently large multiple minus the adjoint is
again nef and big.

On a common smooth resolution, let $P_a$ be $a$ times the pullback of
this semiample adjoint.  The negative-map comparison, using the
negativity lemma \cite[Lemma~3.39]{KollarMori1998}, gives
\[
 aK_W+L\sim_{\Q}P_a+F_a,
 \qquad F_a\ge0
\]
after pullback, with $F_a$ exceptional over the model.  Sections in
divisible degrees are unchanged by this exceptional part, so
$P_a^d=\vol(aK_W+L)$.  All occurrences of $L$ on the resolution mean
its nef pullback.  The Khovanskii--Teissier inequality gives
\[
 (aK_W+L)\cdot L^{d-1}
 \ge P_a\cdot L^{d-1}
 \ge (P_a^d)^{1/d}(L^d)^{(d-1)/d}.
\]
For completeness, the nef mixed-intersection inequalities used here
follow by ample perturbation and passage to the limit from the
surface Hodge index theorem, applied on general complete-intersection
surfaces.  Combining with \eqref{eq:scalar-fibre-hypotheses}, we
obtain the explicit bound
\begin{equation}\label{eq:scalar-volume-bound}
 \vol(aK_W+L)\le (aC'+1)^dL^d,
\end{equation}
which proves the claim.

Pass to a subsequence on which $j=\kappa(W)$ is fixed.
Since $d\le3$, effective Iitaka fibrations in the zero-boundary case
give an integer $m_0=m_0(d)>0$ such that $|m_0K_W|$ gives the Iitaka
fibration; see \cite[Corollary~1.5]{ChenHanLiu2023}.
The case $j=d$ is impossible for large sequence index.  Resolving
that linear system would produce an integral free moving part with
generically finite map and top self-intersection at least one.  Thus
$\vol(K_W)\ge m_0^{-d}$, whereas
$\vol(K_W+L)\to0$ by \eqref{eq:scalar-small-adjoint-volume} and
$L$ has effective multiples.

If $0<j<d$, resolve the Iitaka system, its Stein base, and the induced
map to obtain a morphism of smooth projective varieties
\[
 f:W\longrightarrow T
\]
with connected fibres.  There is an integral big basepoint-free divisor
$H$ on $T$, pulled back from the hyperplane class of the system's image,
such that
\begin{equation}\label{eq:scalar-iitaka-hyperplane}
 f^*H\preceq m_0K_W,
 \qquad H^j\ge1.
\end{equation}
We continue to write $W$ for the resulting smooth model and $L$ for
its pullback.  The previous intersection estimates persist: canonical
divisors under further smooth resolutions acquire effective exceptional
terms, and these pair to zero against $L^{d-1}$.

Let $U$ be a smooth integral geometric generic fibre of $f$ and put
$g=d-j>0$.  Then $K_W|_U\sim K_U$, and $L|_U$ is big because these
fibres move.  We recall explicitly that $\kappa(U)=0$.  A nonzero
canonical multiple on $W$ restricts nontrivially to the generic fibre,
so $\kappa(U)\ge0$.  If a system $|kK_U|$ had positive-dimensional
image, its sections would extend over the generic base after twisting
$kK_W$ by $vf^*H$ for a sufficiently large $v$.  Indeed the coherent
direct image of $kK_W$ has the required generic fibre of sections,
and a sufficiently large multiple of the big divisor $H$ makes it
generically generated, using a decomposition into ample and effective
parts.  By \eqref{eq:scalar-iitaka-hyperplane}, the twist is absorbed
into a further canonical multiple.  The resulting section ratio is
nonconstant on the geometric generic fibre and hence transcendental
over $\C(T)$.  Products with sections of $|m_0K_W|$ also preserve the
$j$ independent base ratios.  This would force $\kappa(W)>j$, a
contradiction.  When $j=0$, simply put $U=W$ and $g=d$.

We next show that
\begin{equation}\label{eq:scalar-small-fibre-volume}
 \vol(K_U+L|_U)\longrightarrow0.
\end{equation}
For $j=0$ this is \eqref{eq:scalar-small-adjoint-volume}.  If $j>0$
and $g=1$, then $\kappa(U)=0$ implies $\deg K_U=0$, and the volume
in question equals $\deg(L|_U)$.  Set $A=f^*H$.  Nef mixed-intersection
log-concavity and \eqref{eq:scalar-iitaka-hyperplane} give
\[
 A^jL^{d-j}
 \le\left(\frac{AL^{d-1}}{L^d}\right)^jL^d
 \le (m_0C')^jL^d\longrightarrow0.
\]
One may add a positive multiple of $L$ to $A$ and pass to the limit
if a mixed intersection vanishes.  The projection formula and
$H^j\ge1$ show that this bounds $\deg(L|_U)$ from above.

The only remaining case is $d=3$, $j=1$, and $g=2$.  Now $T$ is a
smooth curve and $H$ is ample.  A basepoint-free pencil in $|H|$ and
its ramification divisor give $K_T+2H\succeq0$.  Consequently
\begin{equation}\label{eq:scalar-relative-domination}
 K_{W/T}+L+2f^*H\preceq(1+4m_0)K_W+L.
\end{equation}
Choose an effective klt rational boundary linearly equivalent to $L$
from its semiample system.  The twisted weak positivity theorem
\cite[Theorem~1.1]{FujinoWeakPositivity} applies to $f$ with this
boundary.  Thus for sufficiently divisible $l$ the locally free sheaf
\[
 E_l=f_*\OO_W\bigl(l(K_{W/T}+L)\bigr)
\]
is weakly positive.  Here smoothness, projectivity, connected fibres,
and effectivity of the klt boundary are all satisfied; on the smooth
curve a torsion-free direct image is locally free.

Choose $l\ge1$, separately for each member of the sequence, so that
the $l$-th Veronese of the section algebra of $K_U+L|_U$ is generated
in degree one.  Such an $l$ exists by the same big klt adjoint model
and finite-generation theorem used above.  Enlarge it to clear the
chosen rational linear equivalence.  Weak positivity gives $b>0$
for which
\[
 \Sym^bE_l\otimes\OO_T(bH)
\]
is generated by global sections at the generic point.  For every
positive multiple $q$ of $b$, symmetric multiplication and the
degree-one generation on the generic fibre therefore give global
sections of
\[
 ql(K_{W/T}+L)+qf^*H
\]
spanning the degree-$ql$ fibre sections.  Select a basis over
$\C(T)$ among their restrictions, of cardinality
$h^0(U,ql(K_U+L|_U))$.  Multiplying by a basis of $H^0(T,qlH)$
gives linearly independent products over $\C$: any relation can
first be tested using independence over $\C(T)$ and then independence
of the base sections.  Multiplication by a nonzero section of
$(ql-q)H$ places these products in degree $ql$ of
$K_{W/T}+L+2f^*H$.

Curve Riemann--Roch and the asymptotic fibre section formula, with
$l$ fixed and $q\to\infty$ through these multiples, yield
\begin{align}
 \vol\bigl((1+4m_0)K_W+L\bigr)
 &\ge\vol(K_{W/T}+L+2f^*H)\notag\\
 &\ge\frac{d!}{g!}\deg(H)\vol(K_U+L|_U).
 \label{eq:scalar-weak-positive-volume}
\end{align}
The first expression tends to zero by
\eqref{eq:scalar-small-adjoint-volume}, and $\deg H\ge1$.
This proves \eqref{eq:scalar-small-fibre-volume} in every case.

\end{proof}

\subsection{The upper order estimate}

\begin{proof}[Proof of the upper bound in Theorem~\ref{thm:scalar-bounds}]
If no uniform $C$ exists, rational rescaling gives a sequence of the
data in the theorem, with polarizations denoted $M$, such that
\begin{equation}\label{eq:scalar-small-volume}
 M^4\longrightarrow0,
 \qquad \gamma(M;V)>1.
\end{equation}
The proof decreases the dimension while retaining small volume and
an effective canonical upper bound. At each stage we first extract a
moving base component, then pass to a Kodaira-dimension-zero fibre of
its Iitaka fibration, and finally replace a big adjoint by its semiample
model. The process stops when its section-order invariant tends to zero.
Tracing the effective divisor comparisons back will construct covering
families of positive-dimensional proper subvarieties $A\subset V$ whose
geometric generic members satisfy
\begin{equation}\label{eq:scalar-small-zero-members}
 \kappa(A^*)=0,
 \qquad \gamma(M;A)\longrightarrow0.
\end{equation}

All constructions can be made successively on geometric generic
members.  To be precise, choose countable fields of definition for
the sequence.  A geometric generic datum over such a field can be
realized over $\C$ to apply the model and positivity theorems used here.
The resulting finite collection of varieties, morphisms, and divisors
descends again to a field finitely generated over the preceding
generic parameter field.  Spreading it out therefore gives a dominant
extension of the preceding parameter space: the old generic parameter
is never specialized.  Formation of section spaces and of their jet
kernels commutes with field extension, so the generic all-degree order
comparisons are preserved.  There are only finitely many stages for
each variety, and countably many varieties in the sequence.

Begin with a small projective $\Q$-factorialization $Z\to V$, which
exists for the klt variety $V$ by \cite[Corollary~1.4.3]{BCHM2010},
and let $N$ be the pullback of $M$ to $Z$.  We describe a stage of the construction.  After
passing to a subsequence, its dimension $n$ is fixed and the following
properties hold:
\begin{enumerate}[label=(\roman*)]
 \item $Z$ is projective $\Q$-factorial klt and $N$ is big, nef, and
       semiample, with $N^n\to0$;
 \item a smooth projective model of $Z$ has nonnegative Kodaira dimension;
 \item $K_Z\preceq a_0N$ for a fixed positive rational number $a_0$;
 \item $\gamma(N;Z)$ is bounded below by a positive constant.
\end{enumerate}
Initially these assertions follow from \eqref{eq:scalar-small-volume},
canonicality, and crepantness of the small morphism.  If $n=1$, they
cannot all hold, since the order of a section on a smooth curve is
bounded by its degree.  We may therefore assume $n\ge2$.

Choose $n+1$ fixed rational jet levels, equally spaced by $\delta>0$,
below the lower bound in (iv).  For all sufficiently large sequence indices they are above
$(N^n)^{1/n}$.  Apply Lemma~\ref{lem:scalar-tracking}.  Write
$W=S^*$ for a resolution of its geometric generic moving member and
$L=N|_W$, and pass to a subsequence with fixed $d=\dim W<n$.
The divisor $L$ is nef and semiample, and is big because a sweeping
member meets the open set on which the semiample map of $N$ is
generically finite.  Lemma~\ref{lem:scalar-covering} gives
$\kappa(W)\ge0$.  Restricting the effective comparison in (iii) to
the moving member and using the tracking bounds yields
\begin{equation}\label{eq:scalar-small-intersections}
 L^d\longrightarrow0,
 \qquad K_W\cdot L^{d-1}\le C' L^d
\end{equation}
with the fixed constant $C'=a_0+2n/\delta$ at this stage.

Apply Lemma~\ref{lem:scalar-iitaka-volume} to these moving members.
After passing to a subsequence and resolving their Iitaka systems,
let $U$ be the smooth geometric generic Iitaka fibre, with $U=W$ when
$\kappa(W)=0$, and put $g=\dim U$.  Then
\[
 0<g\le d<n,\qquad \kappa(U)=0,\qquad
 \vol(K_U+L|_U)\longrightarrow0.
\]
The restricted divisor $L|_U$ is big, nef, and semiample.  The generic
field and spreading convention above keeps these fibres in covering
families over the preceding parameter space.

To continue the dimension descent, we need a semiample polarization
that still bounds the original restricted sections from above.
Take an effective klt boundary $\Gamma\sim_{\Q}L|_U$ from the free
system, and take a big klt good model of $K_U+\Gamma$, as in the proof
of Lemma~\ref{lem:scalar-iitaka-volume}.
Write $Z'$ for its $\Q$-factorial klt underlying variety and $N'$
for its big semiample adjoint.  Then
\[
 (N')^g=\vol(K_U+L|_U)\longrightarrow0.
\]
On a smooth common resolution with maps $a$ to $U$ and $b$ to $Z'$,
the negative-map comparison reads
\begin{equation}\label{eq:scalar-good-model-comparison}
 a^*(K_U+\Gamma)=b^*N'+F,
 \qquad F\ge0\text{ exceptional over }Z'.
\end{equation}
For any effective $\Delta_U\sim_{\Q}K_U$, we have
\begin{equation}\label{eq:scalar-fixed-part-bound}
 F\le a^*\Delta_U.
\end{equation}
Indeed, for sufficiently divisible $m$, every effective divisor in
$|mb^*N'+mF|$ contains $mF$.  Its pushforward is an effective divisor
in $|mN'|$ on the normal variety $Z'$, and pulling that Cartier divisor
back leaves exactly the exceptional summand $mF$.  But
$ma^*\Delta_U+Y$ belongs to this system when $Y$ is a general member
of the free system $|ma^*(L|_U)|$.  Choose $Y$ avoiding every component
of $F$.  Comparing coefficients proves
\eqref{eq:scalar-fixed-part-bound}.  In particular,
\begin{equation}\label{eq:scalar-polarization-domination}
 a^*(L|_U)\preceq b^*N'.
\end{equation}

If $\gamma(N';Z')\to0$ on a subsequence, we stop.  Otherwise, on a
subsequence it is bounded below by a positive constant, and we repeat
with $(Z',N')$.  The four stage conditions hold: $K_{Z'}\preceq N'$
because the pushed-forward boundary is effective, and a smooth model
of $Z'$ is birational to $U$, which has Kodaira dimension zero.
The dimension has decreased strictly, from $n$ to $g\le d<n$.
The construction must stop, since on a curve maximum normalized
section order is at most the degree.

We verify that the stopping condition proves
\eqref{eq:scalar-small-zero-members} for the original polarization.
Every moving member and fibre used above spreads to a covering
family over the preceding generic parameter space.  Its map to its
image in that space, and ultimately in $V$, is birational: a moving
member meets the open set of isomorphism of each relevant birational
model map.  For a subsequent member in $Z'$, take its strict transform
on the common resolution, so it has an actual morphism back to $U$.
The effective comparisons \eqref{eq:scalar-polarization-domination}
restrict to these members, since they are not contained in the extra
divisors.  Their successive restrictions compare the pullback of the
original $M$ with the final $N'$ on a common resolution.  Multiplication
by the effective differences injects section spaces in divisible
degrees, and at a generic marked point those differences do not
vanish.  Birational pullback between normal projective models preserves
the section spaces.  Thus the final image $A\subset V$ satisfies
\[
 \gamma(M;A)\le\gamma(N';Z')\longrightarrow0,
 \qquad \kappa(A^*)=\kappa(U)=0.
\]
Composing the dominant incidence evaluations proves that these images
form a covering family.  They have positive dimension and are proper
because the first tracking step already decreased dimension.

Mark a generic smooth point on these members and add every $G$-translate
of the marked family.  All data can be defined over one countable
field; invariance of $M$ preserves the order bound.
Lemma~\ref{lem:generic-chain} gives a $G$-equivariant rational quotient,
and Lemma~\ref{lem:chain-order} bounds the section order on its
geometric generic leaf by its dimension times $B$, where $B\to0$.
If the leaf has dimension four, this gives $\gamma(M;V)\le4B$, contradicting
\eqref{eq:scalar-small-volume}.  Otherwise it gives a $G$-equivariant
rational fibration with positive-dimensional base and fibre.  The
geometric generic leaf is of general type by hypothesis.

That leaf is nevertheless covered by members of Kodaira dimension
zero.  To see that the assertion is about generic members, restrict
the marked incidence to its generic leaf name.  A component still
dominates the leaf by marked-point evaluation.  Its generic parameter
is generic in the original parameter space over the original field,
because it is obtained by geometric extension over the generic name.
The quotient is generically constant on each such member, and its
strict transforms give a covering family on a resolution of the leaf.
This contradicts Lemma~\ref{lem:scalar-covering}.  The upper bound
follows.
\end{proof}

\subsection{The lower curve estimate}

\begin{proof}[Proof of the lower bound in Theorem~\ref{thm:scalar-bounds}]
Rescale rationally so that $1\le P^4\le2$.  If no uniform positive
lower bound exists, there is a sequence of marked covering families
of integral curves with
\[
 \frac{P\cdot D}{\mult_xD}\le B,
 \qquad B\longrightarrow0.
\]
Here the marks can be chosen generic over countable fields of
definition.  Indeed, failure of a very general bound still gives a
bad mark after excluding the countably many proper closed subsets
over such a field.  Choose a curve witnessing a strict upper bound
there and spread the curve and its mark over their field of definition.
The mark evaluation is dominant, and degree and multiplicity agree
with those of the geometric generic data.  Include all $G$-translates
if necessary; they obey the same ratio bound.

Take the rational quotient from Lemma~\ref{lem:generic-chain}.
After passing to a subsequence with fixed leaf dimension $h>0$,
Lemma~\ref{lem:chain-order} gives, on its geometric generic leaf $H$,
\begin{equation}\label{eq:scalar-small-leaf}
 P^h\cdot H\le(hB)^h.
\end{equation}
For large sequence index it cannot have dimension four, since
$P^4\ge1$.  Put $s=4-h>0$.  Resolve the leaf fibration to a morphism
$f:X\to T$ between smooth projective varieties with geometrically
integral generic fibre, and let $\pi:X\to V$ be the resulting
birational morphism.  Continue to write $P$ for $\pi^*P$.

Choose a smooth integral very general fibre $F$ of $f$.  It can be
chosen not contained in any of the countably many proper closed sets
excluded by the all-degree upper order bound, nor in the exceptional
locus of $\pi$.  For clarity, over a smooth open of the base, each
proper closed set of $X$ contains the whole fibre only over a proper
closed subset of $T$, by properness and upper semicontinuity of fibre
dimension.  Exclude these countably many subsets first, then choose
a point $x\in F$ outside their intersections with $F$.  Thus $\pi$
is an isomorphism near $x$ and
\begin{equation}\label{eq:scalar-point-upper}
 \gamma(P;V)\le C2^{1/4}
\end{equation}
holds at $\pi(x)$ in every divisible degree.  The degree
$P^h\cdot F$ equals the geometric generic degree in
\eqref{eq:scalar-small-leaf}.

Let $I=I_F\subset\OO_X$.  Since $F$ is a smooth fibre over a smooth
point of the $s$-dimensional base, it is regularly embedded and has
trivial conormal bundle of rank $s$.  Consequently
\[
 I^k/I^{k+1}\simeq
 \Sym^k(\OO_F^{\oplus s})
 \simeq\OO_F^{\oplus\binom{k+s-1}{s-1}}.
\]
The filtration by powers of $I$ therefore bounds the rank of
restriction to the $u$-th thickening, for every $u\ge1$ and every
divisible $l$, by
\begin{equation}\label{eq:scalar-thickening}
 h^0\bigl(X,\OO_X(lP)\otimes(\OO_X/I^u)\bigr)
 \le\binom{u+s-1}{s}\,h^0(F,lP|_F).
\end{equation}
No vanishing theorem uniform in $u$ is required for this inequality.

Fix $D>C2^{1/4}$ and set $u=\lceil Dl\rceil$.  The pullback of $P$
is big and semiample on $X$, and its restriction to $F$ is big and
semiample because $F$ meets the birational isomorphism locus.
The semiample Hilbert-polynomial asymptotics show that the leading
coefficient, after division by $l^4$, on the right-hand side of
\eqref{eq:scalar-thickening} is
\[
 \frac{D^s(P^h\cdot F)}{s!\,h!}
 \le\frac{D^s(hB)^h}{s!\,h!}.
\]
Choose the sequence index first so that this is strictly less than
$1/4!$, hence less than $P^4/4!$.  Then choose a sufficiently large
divisible $l$ for this fixed example.  Since
\[
 h^0(X,lP)=\frac{P^4}{4!}l^4+O(l^3),
\]
restriction to the thickening has a nonzero kernel.  A kernel section
lies in $I^u$, and $I_x^u\subset\mathfrak m_x^u$.  It therefore has
order at least $u\ge Dl$ at $x$.  Normality gives
$\pi_*\OO_X=\OO_V$, so it is a section downstairs as well, with
the same order at the point where $\pi$ is an isomorphism.  This
contradicts \eqref{eq:scalar-point-upper}.  Rescaling back proves the
lower bound, and completes the theorem.
\end{proof}
